%% toc: P1 3 3.7,3.8
\section{The Ultrafilter Lemma and its strength}
\label{sec:p01-w02-p01-009-ultrafilter-lemma}

Definition~\ref{def:ultrafilter-lemma} in \cref{sec:set-prerequisites-06}
states the Ultrafilter Lemma as a principle and records three facts about it:
that it concerns families of subsets of a set, that it belongs to the ambient
set theory rather than to a first-order theory under study, and that over
\(\mathsf{ZF}\) it is strictly weaker than the Axiom of Choice
\cite[ch.~``Sets, Relations, Functions, and Cardinality'', sec.~``Choice
principles used later'', definition ``Ultrafilter Lemma'']{FrozenFormalLogicUltrafilterLemma}.
This section restates exactly that principle, with the same index set, the
same hypotheses and the same conclusion, as a numbered theorem whose standing
is explicit: what proves it, what it is equivalent to over \(\mathsf{ZF}\),
and the published independence result behind the strictness claim.

\emph{Inputs.} A set \(I\) and a family \(\mathcal F\) of subsets of \(I\)
satisfying the four conditions displayed in
Definition~\ref{def:ultrafilter-lemma}; the Axiom of Choice of
\cref{sec:set-prerequisites-06} as an additional hypothesis where stated.
\emph{Conclusions.} A family \(\mathcal U\supseteq\mathcal F\) with the same
four conditions that decides every subset of \(I\) against its complement
(Theorem~\ref{thm:p01-w02-p01-009-ultrafilter-lemma}, part~1); the
\(\mathsf{ZF}\)-provable implication from Choice to the lemma (part~2); the
\(\mathsf{ZF}\)-equivalence of the lemma with the Boolean Prime Ideal
principle and with propositional compactness for arbitrary variable sets
(part~3); and the unprovability of the converse implication, relative to the
consistency of \(\mathsf{ZF}\) (part~4).

\subsection{Assumptions and vocabulary}
\label{sec:p01-w02-p01-009-assumptions}

The ambient theory is \(\mathsf{ZF}\), the classical metatheoretic set theory
fixed at the start of the set chapter; the Axiom of Choice,
\(\mathsf{AC}\), is used only where a statement names it. The Ultrafilter
Lemma is a single sentence of this ambient theory: it quantifies over all
sets \(I\) and all families \(\mathcal F\subseteq\mathcal P(I)\). It is not
a formula of a first-order theory under study, and none of the results below
adds it to an object theory.

For a set \(I\) and a family \(\mathcal F\subseteq\mathcal P(I)\), the
conditions of Definition~\ref{def:ultrafilter-lemma} are
\begin{enumerate}
  \item[(F1)] \(I\in\mathcal F\);
  \item[(F2)] \(\varnothing\notin\mathcal F\);
  \item[(F3)] \(X,Y\in\mathcal F\Rightarrow X\cap Y\in\mathcal F\);
  \item[(F4)] \(X\in\mathcal F\) and \(X\subseteq Y\subseteq I\)
    \(\Rightarrow Y\in\mathcal F\);
\end{enumerate}
and the decision condition on a family \(\mathcal U\subseteq\mathcal P(I)\) is
\begin{enumerate}
  \item[(D)] for every \(X\subseteq I\), exactly one of \(X\) and
    \(I\setminus X\) belongs to \(\mathcal U\).
\end{enumerate}
The definition announces that a later chapter supplies short names for the
two kinds of family. The next lemma identifies them with the vocabulary of
Definition~\ref{def:p01-bpi-filter-ideal}, so that the Boolean results of
\cref{sec:p01-boolean-prime-ideal-and-ultrafilters} apply to them.

\begin{hblemma}{thm:p01-w02-p01-009-terms}{proved=annex:p01-w02-p01-009-terms}{Names for the two families}
Work in \(\mathsf{ZF}\). Let \(I\) be a set and
\(\mathcal F,\mathcal U\subseteq\mathcal P(I)\).
\begin{enumerate}
  \item \(\mathcal F\) satisfies (F1)--(F4) iff \(\mathcal F\) is a proper
    set filter on \(I\), that is, a proper Boolean filter of the algebra
    \(\mathcal P(I)\) with meet \(\cap\), join \(\cup\), complement
    \(I\setminus X\), bottom \(\varnothing\) and top \(I\).
  \item \(\mathcal U\) satisfies (F1)--(F4) and (D) iff \(\mathcal U\) is an
    ultrafilter on \(I\), that is, a proper set filter on \(I\) that is
    maximal under inclusion among proper set filters on \(I\).
  \item If \(I=\varnothing\), no family satisfies (F1)--(F4).
\end{enumerate}
No choice principle is used.
\end{hblemma}

\subsection{The statement and its standing}
\label{sec:p01-w02-p01-009-statement}

Part~1 below carries the principle of Definition~\ref{def:ultrafilter-lemma}
with its notation, hypotheses and conclusion unchanged; parts~2--4 make its strength over \(\mathsf{ZF}\)
precise. ``Strictly weaker over \(\mathsf{ZF}\) than \(\mathsf{AC}\)'' means
the conjunction of part~2 with part~4.

\begin{hbtheorem}{thm:p01-w02-p01-009-ultrafilter-lemma}{proved=annex:p01-w02-p01-009-lemma}{Ultrafilter Lemma}
\begin{enumerate}
  \item \emph{(The lemma.)} Assume \(\mathsf{ZF}+\mathsf{AC}\). Let \(I\) be a
    set and let \(\mathcal F\subseteq\mathcal P(I)\) satisfy
    \[
      I\in\mathcal F,\quad\varnothing\notin\mathcal F,\quad
      X,Y\in\mathcal F\Rightarrow X\cap Y\in\mathcal F,
    \]
    and
    \[
      X\in\mathcal F\text{ and }X\subseteq Y\subseteq I
      \Rightarrow Y\in\mathcal F.
    \]
    Then there is \(\mathcal U\subseteq\mathcal P(I)\) containing
    \(\mathcal F\) that satisfies the same conditions and, for every
    \(X\subseteq I\), contains exactly one of \(X\) and \(I\setminus X\).
  \item \emph{(Choice suffices.)} \(\mathsf{ZF}\) proves the implication from
    \(\mathsf{AC}\) to the principle of part~1, stated for all \(I\) and
    \(\mathcal F\).
  \item \emph{(Equivalents.)} Over \(\mathsf{ZF}\), that principle is
    equivalent to \(\mathsf{BPI}\), to Boolean ideal extension and to Boolean
    filter extension (Definition~\ref{def:p01-bpi-principles}), and to
    propositional compactness for arbitrary sets of variables
    (Definition~\ref{def:p01-bpi-propositional-compactness}).
  \item \emph{(Choice is not implied.)} If \(\mathsf{ZF}\) is consistent, then
    \(\mathsf{ZF}\) does not prove the implication from that principle to
    \(\mathsf{AC}\). This part is derived from the theorem of Halpern and
    L\'evy that consistency of \(\mathsf{ZF}\) implies consistency of
    \(\mathsf{ZF}+\mathsf{BPI}+\neg\mathsf{AC}\) \cite{HalpernLevy1971},
    recorded as Theorem~\ref{thm:p01-bpi-choice-separation}.
\end{enumerate}
Consequently the Ultrafilter Lemma is strictly weaker over \(\mathsf{ZF}\)
than the Axiom of Choice, relative to the consistency of \(\mathsf{ZF}\), and
it is not a second name for that axiom. It is a metatheoretic choice
principle about families of subsets, not a formula of a first-order theory
under study.
\end{hbtheorem}

The proofs of parts~1--3 are complete in the annex and use only results
proved in this handbook. Part~4 is derived there in full from one external
premise, the Halpern--L\'evy relative-consistency theorem, whose model
construction is not reproduced in this handbook. Repick\'y recalls that
theorem as the existence of a transitive model of \(\mathsf{ZF}\) in which the
Boolean Prime Ideal Theorem holds and Choice fails
\cite[p.~543, Theorem~1]{Repicky2015BPI}. Starting from a ground universe
assumed to satisfy \(\mathsf{ZFC}\), he proves that the prime-ideal principle
holds in such a model \cite[pp.~543--546]{Repicky2015BPI}.

\subsection{Applying the lemma: side conditions}
\label{sec:p01-w02-p01-009-side-conditions}

A family to be extended is often given by a few sets rather than as a
family satisfying (F1)--(F4). The following lemma states exactly when such a
family lies inside a deciding family, and isolates the single step that needs
the lemma.

\begin{hblemma}{thm:p01-w02-p01-009-side-conditions}{proved=annex:p01-w02-p01-009-side-conditions}{When a family of subsets can be extended}
Let \(I\) be a set and \(\mathcal E\subseteq\mathcal P(I)\). Say that
\(\mathcal E\) has the finite intersection property iff
\(X_1\cap\cdots\cap X_n\ne\varnothing\) for every \(n\in\mathbb N\) and all
\(X_1,\ldots,X_n\in\mathcal E\), the intersection of the empty list (\(n=0\))
being \(I\). Put
\[
  \begin{aligned}
  \mathcal E^{\uparrow}=\{\,Y\subseteq I:{}&X_1\cap\cdots\cap X_n\subseteq Y
  \text{ for some }n\in\mathbb N\\
  &\text{and some }X_1,\ldots,X_n\in\mathcal E\,\}.
  \end{aligned}
\]
\begin{enumerate}
  \item (\(\mathsf{ZF}\).) If some \(\mathcal U\supseteq\mathcal E\) satisfies
    (F1)--(F4), then \(\mathcal E\) has the finite intersection property.
  \item (\(\mathsf{ZF}\).) If \(\mathcal E\) has the finite intersection
    property, then \(\mathcal E^{\uparrow}\) contains \(\mathcal E\),
    satisfies (F1)--(F4), and is contained in every family that contains
    \(\mathcal E\) and satisfies (F1)--(F4).
  \item (\(\mathsf{ZF}\) together with the principle of
    Theorem~\ref{thm:p01-w02-p01-009-ultrafilter-lemma}, part~1.) Some
    \(\mathcal U\supseteq\mathcal E\) satisfies (F1)--(F4) and (D) iff
    \(\mathcal E\) has the finite intersection property.
\end{enumerate}
\end{hblemma}

\begin{hbmethod}{def:p01-w02-p01-009-procedure}{Using the lemma on a given family}
Given \(I\) and \(\mathcal E\subseteq\mathcal P(I)\):
\begin{enumerate}
  \item Check the finite intersection property. If it fails, stop: by
    Lemma~\ref{thm:p01-w02-p01-009-side-conditions}, part~1, no deciding
    extension exists, whatever choice principle is assumed.
  \item Otherwise form \(\mathcal E^{\uparrow}\); it satisfies (F1)--(F4)
    without any choice principle.
  \item Apply the lemma once to \(\mathcal E^{\uparrow}\). The step requires
    the principle of Theorem~\ref{thm:p01-w02-p01-009-ultrafilter-lemma},
    part~1, or one of its equivalents in part~3. When \(I\) is finite it
    requires no choice principle, by Corollary~\ref{thm:p01-bpi-set-boundaries}.
\end{enumerate}
Side conditions: the procedure yields existence only. It names no particular
extension, and different extensions may exist
(Example~\ref{def:p01-bpi-worked-finite}).
\end{hbmethod}

\begin{hbexample}{def:p01-w02-p01-009-worked}{One family that extends uniquely and one that cannot extend}
Let \(I=\{1,2,3\}\).

\emph{A family that extends.} Take \(\mathcal E=\{\{1,2\},\{2,3\}\}\). The
intersections of its finite lists are \(I\), \(\{1,2\}\), \(\{2,3\}\) and
\(\{1,2\}\cap\{2,3\}=\{2\}\), all nonempty, so the finite intersection
property holds. The subsets of \(I\) containing one of these are exactly those
containing \(2\), so
\(\mathcal E^{\uparrow}=\{\{2\},\{1,2\},\{2,3\},\{1,2,3\}\}\). This family
already satisfies (D): for each \(X\subseteq I\), \(2\) lies in exactly one of
\(X\) and \(I\setminus X\). It is also the only deciding extension: if a
family \(\mathcal U\supseteq\mathcal E^{\uparrow}\) satisfying (F1)--(F4)
contained some \(X\) with \(2\notin X\), it would contain \(I\setminus X\),
which contains \(2\), and hence \(X\cap(I\setminus X)=\varnothing\), against
(F2). No choice principle is involved, because \(I\) is finite and the
extension is written down.

\emph{A family that cannot extend.} Take \(\mathcal E'=\{\{1,2\},\{3\}\}\).
Here \(\{1,2\}\cap\{3\}=\varnothing\), so the finite intersection property
fails. By Lemma~\ref{thm:p01-w02-p01-009-side-conditions}, part~1, no family
satisfying (F1)--(F4) contains \(\mathcal E'\); in particular the Ultrafilter
Lemma has nothing to extend, and neither does \(\mathsf{AC}\).
\end{hbexample}

\begin{hbwarning}{def:p01-w02-p01-009-limits}{What the strength claim does and does not say}
\begin{enumerate}
  \item The strictness claim is relative. Part~4 of
    Theorem~\ref{thm:p01-w02-p01-009-ultrafilter-lemma} assumes that
    \(\mathsf{ZF}\) is consistent; an inconsistent theory proves every
    implication. It is a statement about provability in \(\mathsf{ZF}\), not
    an assertion that Choice is false.
  \item Part~4 rests on a cited construction. The handbook derives part~4 from
    the Halpern--L\'evy theorem \cite{HalpernLevy1971} but does not construct
    the model in which the prime-ideal principle holds and Choice fails.
  \item Nothing here asserts that the Ultrafilter Lemma itself is unprovable
    in \(\mathsf{ZF}\). Part~4 compares the lemma with Choice; it gives no
    information on whether \(\mathsf{ZF}\) alone proves the lemma.
  \item The lemma asserts existence and supplies no description of the
    extension. For an infinite set, every extension of the cofinite filter
    omits all finite sets (Corollary~\ref{thm:p01-bpi-set-boundaries}), but
    none of the results here names such an extension.
  \item An argument that applies the lemma once and then also selects one
    witness for each member of an index set uses more than the lemma; the
    model-theoretic uses are mapped in
    Boundary~\ref{def:p03-ultrafilter-choice-boundary}.
  \item Part~3 concerns propositional compactness for an arbitrary set of
    variables. Compactness for a fixed enumerable variable set has a separate
    proof without any choice principle and is not equivalent to the lemma by
    the results of this section
    (Theorem~\ref{thm:p01-bpi-propositional-compactness}).
\end{enumerate}
\end{hbwarning}

\subsection{Relationships}
\label{sec:p01-w02-p01-009-relationships}

The theorem is the statused form of Definition~\ref{def:ultrafilter-lemma}
and is used wherever that definition is cited. The Boolean development it
relies on is \cref{sec:p01-boolean-prime-ideal-and-ultrafilters}: the
filter-ideal calculus (Lemma~\ref{thm:p01-bpi-duality}), the four-way
equivalence (Theorem~\ref{thm:p01-bpi-equivalence}), the compactness bridge
(Theorem~\ref{thm:p01-bpi-propositional-compactness}, restated in Part~2 as
Theorem~\ref{thm:p02-arbitrary-prop-compactness-bpi}), the Zorn application
(Theorem~\ref{thm:p01-bpi-ac-implication}) and the cited separation
(Theorem~\ref{thm:p01-bpi-choice-separation}). The Zorn step itself is
Theorem~\ref{thm:p06-ac-to-zorn}. The placement of the prime-ideal principle
among other weak choice principles is Theorem~\ref{thm:p06-choice-implications}
in \cref{sec:p06-cardinals-choice}, which also records the Halpern--L\'evy
separation.

\subsection{Sources}
\label{sec:p01-w02-p01-009-sources}

The statement of part~1, its two displayed condition blocks and its
qualifications follow the inherited manuscript's definition
\cite[ch.~``Sets, Relations, Functions, and Cardinality'', sec.~``Choice
principles used later'', definition ``Ultrafilter Lemma'']{FrozenFormalLogicUltrafilterLemma}.
That definition asserts the strictness claim without a reference and without
the consistency hypothesis; here the claim carries both. The independence
result is that of Halpern and L\'evy \cite{HalpernLevy1971}, as recalled by
Repick\'y \cite[p.~543, Theorem~1]{Repicky2015BPI}. The proofs of
parts~1--4 and of the two lemmas are editorial arguments written for this
handbook.
